% ===========================================================================
\section{Convex and smooth: guaranteed descent}\label{sec:smooth}
\warmth{0}\quad\whisper{Convexity rules out traps; smoothness caps the curvature; together they force $O(1/k)$.}

\begin{strand}
Two assumptions make the toy rigorous in any dimension. \keyword{Convexity} says the bowl has no
false bottoms, every stationary point is a global minimum. \keyword{$L$-smoothness} says the
gradient never changes faster than rate $L$, so a step of size $1/L$ is always safe. Together they
guarantee the value falls every step, and reaches the minimum at rate $O(1/k)$.
\end{strand}

\block{The two assumptions}

\begin{definition}[convex, $L$-smooth]\label{def:smooth}
$f$ is \keyword{convex} if it lies above each tangent: $f(y)\ge f(x)+\inner{\nabla f(x)}{y-x}$. It
is \keyword{$L$-smooth} if $\nabla f$ is $L$-Lipschitz, equivalently it lies below each tangent
parabola:
\[
  f(y)\;\le\; f(x)+\inner{\nabla f(x)}{y-x}+\frac{L}{2}\norm{y-x}^2 .
\]
\end{definition}

\block{Descent, then the rate}

\begin{lemma}[descent lemma]\label{lem:descent}
With step $\eta=1/L$, a gradient step strictly decreases $f$:
\[
  f(x_{k+1})\;\le\; f(x_k)-\frac{1}{2L}\,\norm{\nabla f(x_k)}^2 .
\]
\end{lemma}
\begin{proof}
Put $y=x_{k+1}=x_k-\tfrac1L\nabla f(x_k)$ into the smoothness inequality. \TODO{expand the two terms; the cross term $-\tfrac1L\norm{\nabla f}^2$ and $+\tfrac{1}{2L}\norm{\nabla f}^2$ combine to $-\tfrac1{2L}\norm{\nabla f}^2$}
\end{proof}

\begin{theorem}[$O(1/k)$ for convex, smooth]\label{thm:smoothrate}
If $f$ is convex and $L$-smooth, then gradient descent with $\eta=1/L$ satisfies
\[
  f(x_k)-f^\star\;\le\;\frac{L\,\norm{x_0-\xstar}^2}{2k} .
\]
\end{theorem}
\begin{proof}
Skeleton: convexity bounds the gap $f(x_k)-f^\star\le\inner{\nabla f(x_k)}{x_k-\xstar}$; combine with
the descent lemma and telescope the distances $\norm{x_k-\xstar}^2$. \TODO{carry out the telescoping sum and divide by $k$}
\end{proof}

\begin{yourturn}
Read the rate as a budget. To reach $f(x_k)-f^\star\le\varepsilon$, Theorem~\ref{thm:smoothrate}
needs
\[
  k\;=\;O\!\big(\fillin[1.5cm]\big)\ \text{steps.}
\]
Halving $\varepsilon$ \fillin[2.5cm] the number of steps. \recall{$k=O(1/\varepsilon)$; halving $\varepsilon$ \emph{doubles} the steps. This ``sublinear'' rate is honest but slow, which is exactly what strong convexity will fix.}
\end{yourturn}

\begin{strand}
A note on $L$: it is the largest eigenvalue of the Hessian (for a quadratic, $L=\lambda_{\max}$), so
the safe step $1/L$ is set by the \emph{steepest} curvature. If other directions are much flatter,
that one step size wastes their progress, exactly the complaint §\ref{sec:cond} makes precise.
\end{strand}
